% BEGIN PROFESSOR FEEDBACK PF33: Expanded Related Work
\chapter{Related Work}
\label{chap:relatedwork}

Research on malicious blockchain activity differs in what is represented and
what is predicted. Some methods classify addresses from transaction summaries,
while others learn from their neighbourhoods or from ordered sequences of
transfers. The datasets also differ in whether labels describe transactions,
addresses, trust relations, or particular types of abuse. This chapter reviews
these approaches and explains how the present evaluation relates to them.

\section{Graph-Based Detection of Malicious Behaviour}
\label{sec:related_graph_detection}

A transaction graph provides more context than a list of independent payments.
However, the choice of graph determines which interactions are visible.
\citet{lee2020ethereum} study several Ethereum interaction networks, covering
ordinary transactions, contract interactions, and token transfers. Their
comparison shows why the graph representation must be specified before its
properties can be interpreted. This is relevant here because an ordinary
transaction graph and a graph that also includes token transfers do not describe
the same set of activities.

Early supervised detection methods use address-level summaries of such
networks. \citet{ostapowicz2019detecting} compare random forests, support vector
machines (SVMs), and Extreme Gradient Boosting (XGBoost) on more than 300,000
Ethereum accounts. They also examine how individual features affect the
classification. Their work establishes a practical starting point for the
present thesis: conventional address features already contain useful detection
information, so any additional representation needs to be tested against them.

\citet{chen2022phishing} construct 149 features from an address's direct
transactions, drawing on patterns associated with fraudulent bank accounts.
Their method filters the features and uses a random forest for phishing
detection, reporting 98.89\% precision on their dataset. The features have a
clear connection to transaction behaviour, but they summarise the target's
activity rather than count ordered interactions among several addresses.

\citet{sanjalawe2023atd} also derive features from a weighted directed
multigraph. Their approach combines feature selection with a semi-supervised
generative adversarial network and reports 95.06\% detection accuracy.
The selected features include structural measures and summaries of transaction
timing. This is different from representing the exact order of transfers in a
small connected subgraph. The reported scores from these studies cannot be
compared directly with each other or with the present results, because the
datasets, class proportions, and evaluation procedures differ.

Together, these studies motivate the conventional feature baseline in this
thesis. Transaction counts, values, activity durations, and static graph
properties can identify differences between malicious and benign addresses.
The purpose of adding local motif-role counts is to test whether short
transaction sequences contribute information beyond these summaries.

\section{Network Embeddings and Time-Aware Learning}
\label{sec:related_temporal_learning}

Network embeddings replace a manually designed feature vector with a learned
representation of an address and its neighbourhood.
\citet{wu2020phishers} introduce trans2vec for Ethereum phishing detection.
The method incorporates transaction amounts and timestamps into biased random
walks and then applies a one-class support vector machine. It therefore already
uses temporal information. Its representation is learned from walks, however,
rather than expressed as separate counts of predefined local roles.

Time-aware graph modelling is also studied outside fraud classification.
\citet{wu2020temporal} represent Ethereum records as a temporal weighted
multidigraph and compare walk strategies that account for transaction times and
amounts. Their evaluation concerns temporal link prediction. This supports the
importance of retaining repeated and timestamped interactions, but link
prediction is a different task from deciding whether an address is malicious.

Graph neural networks provide another way to combine an object's features with
information from its neighbours. In the Bitcoin setting,
\citet{weber2019bitcoin} compare graph convolutional networks with conventional
classifiers on the Elliptic dataset. Random forests perform strongly in their
experiments. The result is a useful reminder that a graph-based model does not
automatically improve on a feature-based classifier, even when the underlying
data is relational.

More recent Ethereum work models changing neighbourhoods explicitly.
\citet{ghoshcatalog2025} combine local and global temporal dependencies for
phishing detection. Their model processes sequences of transaction
neighbourhoods and learns how behaviour changes over time. This differs from a
fixed catalogue of short event patterns, but it rules out a general claim that
earlier blockchain detection methods ignore time.

The distinction relevant to this thesis is therefore between representations,
not simply between methods that use time and methods that do not. Temporal
motif fingerprints provide explicit counts for particular ordered patterns.
Embeddings and neural models learn their representations from the data. The
present experiments compare fingerprints with conventional features. They do
not benchmark all of these learned alternatives or isolate the causal
contribution of temporal ordering.

\section{Datasets for Blockchain Transaction Analysis}
\label{sec:relwork_datasets}

A useful detection dataset needs both behavioural observations and labels for
the objects being classified. It also needs sufficient transaction detail for
the proposed representation. An address-level table of summary features cannot
by itself reconstruct the order of interactions between several addresses.

\citet{alemari2021labeled} introduce the Benchmark Labelled Transactions of
Ethereum (BLTE) dataset as a public resource for evaluating detection methods.
The later study by \citet{sanjalawe2023atd} describes its source as Ethereum
Classic and uses abnormal-address information from EtherScamDB. BLTE is relevant
as an established transaction-based benchmark. Its source coverage and binary
labels must nevertheless be distinguished from those of a mainnet Ethereum
address-classification task.

The Elliptic dataset released by \citet{weber2019bitcoin} contains more than
200,000 Bitcoin transactions connected by directed payment flows. Its nodes
represent transactions, not persistent Ethereum addresses. It supports the
study of illicit transaction detection over time, but its graph construction
and prediction unit differ from those used here.

The Real-world dataset of Cryptocurrency Addresses with Transaction Profiles
(Real-CATS), introduced by \citet{shi2025realcats}, combines 103,203 criminal
addresses and 106,196 benign addresses across Bitcoin and Ethereum. Criminal
labels are drawn from abuse reports, while benign addresses are collected from
exchange customers. Alongside address profiles, the release contains
transaction histories that support alternative feature extraction. Its
Ethereum categories make it useful for the category comparisons in this
thesis. These labels describe reported abuse, however, and do not establish
that every address in a category belongs to one coordinated campaign.
A supplementary dataset of cryptocurrency addresses and transaction profiles,
named Sup-CATS, provides a separate deployment-oriented evaluation sample.

The Ethereum Fraud Dataset by Activity provides another source of transaction
histories and address labels \citep{fesevu2025ethfraud}. The release includes raw
transfers with sender, receiver, timestamp, and transaction metadata, together
with labels assembled from Etherscan and curated sources. This makes it
possible to compute conventional features and temporal fingerprints from the
same underlying observations. Its role in the present study is as a data
resource, not as evidence for the performance of a particular detection method.

These resources are complementary rather than interchangeable. Their labels
refer to different objects, and their transaction histories cover different
parts of the underlying networks. Chapter~\ref{ch:method} therefore separates
the role of each source from the selection and graph-construction rules used
in the experiments. Comparisons between feature sets use the same addresses
and data splits, so that differences in their scores are not caused by changing
the evaluation sample.

\section{Temporal Motifs and Local Role Representations}
\label{sec:related_temporal_motifs}

Motif analysis begins with the idea that small patterns can reveal structure
that is not captured by single-node statistics.
\citet{milo2002network} identify network motifs through overrepresentation
relative to randomised networks. Counting a candidate pattern and establishing
that it is statistically overrepresented are therefore separate steps.
\citet{McDonnell2014} make a further distinction between a global subgraph
count and a node's participation in its local roles. Their motif-role
fingerprints describe the positions occupied by individual nodes in directed
patterns. This provides the basis for an address-level representation rather
than a single summary of the entire transaction graph.

\citet{kovanen2011temporal} extend pattern analysis to event sequences whose
similarity depends on both topology and temporal order.
\citet{paranjape2017motifs} develop an efficient counting framework for
$\delta$-temporal motifs and show that pattern counts and relevant time scales
differ across networks. In this framework, an ordered sequence must fit within
a bounded time span. The present thesis uses this counting definition while
referring to untested patterns as candidate motifs. Their predictive value is
evaluated without claiming network overrepresentation.

Several studies connect these ideas to financial behaviour.
\citet{lal2021money} examine temporal cycles in Ethereum money trails and find
different cyclic patterns associated with suspicious activities. Their
heuristic analysis concerns paths through which funds move. It supports
examining connected transfer sequences, but it is not the same evaluation as
testing a fixed local-role fingerprint against a conventional feature baseline.

\citet{wu2022mixing} combine temporal and attributed heterogeneous patterns
with other transaction features to detect Bitcoin mixing services. They use
positive and unlabelled learning because reliable negative labels are limited.
This is direct evidence that motif-based features have already been used in
supervised blockchain analysis. The target is membership in a mixing service,
which is different from the binary and category labels considered here.

There is also a direct Ethereum precedent.
\citet{wang2023temporal} study phishing detection using temporal motif
features extracted from transaction subgraphs. The combination of blockchain
data, labelled phishing addresses, and temporal motif features is therefore
not new. The present work builds on this line of research through its
evaluation of local-role fingerprints, temporal-window sensitivity, and
differences between recorded behaviour categories.

\citet{arnold2024motifs} analyse three-transfer patterns on two Bitcoin
datasets and one non-fungible token (NFT) transaction dataset. Their results
show that global counts can be dominated by a small number of participants or
short periods of activity. Local signatures and pattern completion times
reveal information that these global totals hide. They also examine randomised
timestamp references and the ambiguity of ordering transactions within a block.
These observations motivate both the address-level fingerprints and the
equal-timestamp sensitivity checks in this thesis. Their analysis is primarily
descriptive, rather than a held-out comparison of malicious-address classifiers.

Outside blockchain networks, \citet{liu2023temporal} evaluate temporal motif
features for fraud detection on Mercari transactions and a synthetic
J.P.\ Morgan Chase network. They report improvements over simple graph
features and the embedding methods Large-scale Information Network Embedding
(LINE) and node2vec. Their Venmo analysis addresses different prediction and
description tasks. These findings support the usefulness of temporal patterns
in financial networks, but do not establish that the same improvement will
hold for Ethereum addresses.

A different use of motifs appears in \citet{mo2022motifgcn}.
The authors combine motif matrices and temporal graph convolutional learning
on the Bitcoin-Alpha and Bitcoin over-the-counter (Bitcoin-OTC) signed trust
networks. Their revised model predicts the signs of trust links. The motifs
describe structure within snapshots, while temporal information is combined
across snapshots. These are neither Ethereum payment histories nor direct
counts of ordered three-transfer occurrences, so the task and representation
need to be kept separate from the present approach.

\section{Position of This Thesis}
\label{sec:related_position}

The reviewed literature already contains time-aware embeddings, temporal graph
models, and motif-based blockchain detection. The contribution of this thesis
is an evaluation of a particular local motif-role representation under a common
protocol. It compares conventional features, fingerprints alone, and their
combination on the same held-out addresses. It also examines how the permitted
time span affects the representation, whether recorded abuse categories are
separable through fingerprints, and whether category-specific windows improve
on a shared window.

These comparisons address the four research questions introduced in
Chapter~\ref{ch:introduction}. The evaluation also records how the results
change when transfers with equal timestamps are ordered differently. Its scope
is the predictive usefulness and stability of the fingerprints under the
available labels and graph coverage. It does not establish that temporal order
alone causes an improvement or that the counted candidates are statistically
overrepresented network motifs.
% END PROFESSOR FEEDBACK PF33
